\begin{frame}{이진 GDA의 모델: 공유 공분산}
  \[ y \sim \mathrm{Bernoulli}(\phi), \hspace{2em}
     x\mid y{=}0 \sim \mathcal{N}(\mu_0,\Sigma), \hspace{2em}
     x\mid y{=}1 \sim \mathcal{N}(\mu_1,\Sigma) \]
  \vskip0.4em
  \begin{itemize}
    \item 두 클래스가 \textbf{같은 공분산 행렬 $\Sigma$ 를 공유}하되
      평균 $\mu_0,\mu_1$ 만 다름.
    \item 직관: 같은 ``모양''으로 퍼져 있고 \textbf{중심 위치만 다른}
      두 구름.
    \item 파라미터 $\phi,\mu_0,\mu_1,\Sigma$ 는 학습 데이터의
      \textbf{평균·공분산을 그대로 계산}해서 구함 (최대우도추정,
      MLE).
    \item \textbf{경사하강법 없이 닫힌 형태}로 바로 구해짐 ---
      로지스틱회귀와의 실전적 차이.
  \end{itemize}
\end{frame}
